\documentclass[a4paper]{article}

% \usepackage[default]{fontsetup}

\usepackage{fancyhdr}
\usepackage{extramarks}
\usepackage{amsmath}
\usepackage{amsthm}
\usepackage{amsfonts}
\usepackage{tikz}
\usepackage[plain]{algorithm}
\usepackage{algpseudocode}
\usepackage{enumerate}

\usetikzlibrary{automata,positioning}

%
% Basic Document Settings
%  

\topmargin=-0.2in
\evensidemargin=0in
\oddsidemargin=0in
\textwidth=6.5in
\textheight=9.5in
\headsep=0.25in

\linespread{1.1}

\pagestyle{fancy}
\lhead{\hmwkAuthorName}
\chead{\hmwkClass : \hmwkTitle}
\rhead{\firstxmark}
\lfoot{\lastxmark}
\cfoot{\thepage}

\renewcommand\headrulewidth{0.4pt}
\renewcommand\footrulewidth{0.4pt}

\setlength\parindent{0pt}

%
% Create Problem Sections
%

\newcommand{\enterProblemHeader}[1]{
    \nobreak\extramarks{}{Problem \arabic{#1} continued on next page\ldots}\nobreak{}
    \nobreak\extramarks{Problem \arabic{#1} (continued)}{Problem \arabic{#1} continued on next page\ldots}\nobreak{}
}

\newcommand{\exitProblemHeader}[1]{
    \nobreak\extramarks{Problem \arabic{#1} (continued)}{Problem \arabic{#1} continued on next page\ldots}\nobreak{}
    \stepcounter{#1}
    \nobreak\extramarks{Problem \arabic{#1}}{}\nobreak{}
}

\newcommand*\circled[1]{\tikz[baseline=(char.base)]{
		\node[shape=circle,draw,inner sep=2pt] (char) {#1};}}


\setcounter{secnumdepth}{0}
\newcounter{partCounter}
\newcounter{homeworkProblemCounter}
\setcounter{homeworkProblemCounter}{1}
\nobreak\extramarks{Problem \arabic{homeworkProblemCounter}}{}\nobreak{}

%
% Homework Problem Environment
%
% This environment takes an optional argument. When given, it will adjust the
% problem counter. This is useful for when the problems given for your
% assignment aren't sequential. See the last 3 problems of this template for an
% example.
%

\newenvironment{homeworkProblem}[2][-1]{%
    \ifnum#1>0\relax
        \setcounter{homeworkProblemCounter}{#1}
    \fi
    \section{Problem \arabic{homeworkProblemCounter} (#2 points)}
    \setcounter{partCounter}{1}
    \enterProblemHeader{homeworkProblemCounter}
}{%
    \exitProblemHeader{homeworkProblemCounter}
}

%
% Homework Details
%   - Title
%   - Class
%   - Due date
%   - Name
%   - Student ID

\newcommand{\hmwkTitle}{Homework\ \#06}
\newcommand{\hmwkClass}{Probability \& Statistics for EECS}
\newcommand{\hmwkDueDate}{2026-11-1}
\newcommand{\hmwkAuthorName}{Jianxuan Yao}
\newcommand{\hmwkAuthorID}{2025521024}


%
% Title Page
%

\title{
    \vspace{2in}
    \textmd{\textbf{\hmwkClass:\\  \hmwkTitle}}\\
    \normalsize\vspace{0.1in}\small{Due\ on\ \hmwkDueDate\ at 23:59}\\
	\vspace{4in}
}

\author{
	Name: \textbf{\hmwkAuthorName} \\
	Student ID: \hmwkAuthorID}
\date{}

\renewcommand{\part}[1]{\textbf{\large Part \Alph{partCounter}}\stepcounter{partCounter}\\}

%
% Various Helper Commands
%

% Useful for algorithms
\newcommand{\alg}[1]{\textsc{\bfseries \footnotesize #1}}
% For derivatives
\newcommand{\deriv}[1]{\frac{\mathrm{d}}{\mathrm{d}x} (#1)}
% For partial derivatives
\newcommand{\pderiv}[2]{\frac{\partial}{\partial #1} (#2)}
% Integral dx
\newcommand{\dx}{\mathrm{d}x}
% Alias for the Solution section header
\newcommand{\solution}{\textbf{\large Solution}}
% Probability commands: Expectation, Variance, Covariance, Bias
\newcommand{\E}{\mathrm{E}}
\newcommand{\Var}{\mathrm{Var}}
\newcommand{\Cov}{\mathrm{Cov}}
\newcommand{\Bias}{\mathrm{Bias}}
\begin{document}


% \maketitle
% \thispagestyle{empty}
% \pagebreak

\date{
Due on Nov. 1, 2026, 23:59 UTC+8}
\title{SI 140A.02  Probability \& Statistics for EECS, Fall 2026 \\
Homework 6}
\maketitle
Read all the instructions below carefully before you start working on the assignment, and before you make a submission.
\begin{itemize}
    \item You are required to submit your homework in PDF format generated by \textbf{LaTeX}, otherwise you will get no points of this homework.
    \item You are required to write down all the major steps towards making your conclusions; otherwise you may obtain limited points of the problem.
    \item Write your homework in English; otherwise you will get no points of this homework.
    \item Any form of plagiarism will lead to $0$ point of this homework. 
\end{itemize}
\newpage

\begin{homeworkProblem}{20}

Let $n\ge 3$, and let $X_1,\ldots,X_n$ be binary random variables
with $\Pr(X_i=1)=p_i\in(0,1)$. Suppose that every collection of
$n-1$ variables is mutually independent. No symmetry of the
joint distribution is assumed.

\begin{enumerate}
\item[(a)] Determine all possible joint PMFs and the exact
admissible range of their free parameter. Prove both necessity
and sufficiency. 
\item[(b)] Suppose additionally that $n=4$, $p_i=1/2$ for all $i$,
and all $16$ joint outcomes have strictly positive probability.
Determine the exact range of
$\Pr(X_1+X_2+X_3+X_4=2)$.
Does the value $3/8$ force mutual independence?
\end{enumerate}
\end{homeworkProblem}

\solution

\begin{enumerate}
    \item[(a)]
    
    Let $X = \left\{X_1, X_2, \dots, X_n\right\}$ be the set of random variables and $x = \left(x_1, x_2, \dots, x_n\right)$ be a realization of these variables.

    Let the independent distribution be $P_0(x) = P(X_1 = x_1, X_2 = x_2, \dots, X_n = x_n) = \prod_{i=1}^n p_i^{x_i}(1-p_i)^{1-x_i}$
    and the actual distribution be $P(x) = P(X_1 = x_1, X_2 = x_2, \dots, X_n = x_n)$.

    Let $\Delta(x) = P(x) - P_0(x)$.

    For any given $x_1, \dots, x_{k-1}, x_{k+1}, \dots, x_n$, we have

    $$
    \sum_{x_k \in \{0, 1\}} P(x) = \prod_{j \ne k} p_j^{x_j} (1 - p_j)^{1 - x_j} = \sum_{x_k \in \{0, 1\}} P_0(x)
    $$

    $$
    \Delta(x_1, \dots, x_{k-1}, 1, x_{k+1}, \dots, x_n) + \Delta(x_1, \dots, x_{k-1}, 0, x_{k+1}, \dots, x_n) = 0
    $$

    Let $\theta = \Delta(1, 1, \dots, 1)$, we have
    
    $$
    P(x) = P_0(x) + \theta \cdot (-1)^{n - \sum_{i=1}^n x_i}
    $$

    $$
    P(x) \ge 0 \iff P_0(x) + \theta (-1)^{n - \sum x_i} \ge 0
    $$

    When $n - \sum x_i$ is even, $(-1)^{n - \sum x_i} = 1$. We need $\theta \ge - P_0(x)$, so $\theta \ge - \min_{\substack{x \in \{0, 1\}^n \\ n - \sum x_i \text{ is even}}} P_0(x)$.

    When $n - \sum x_i$ is odd, $(-1)^{n - \sum x_i} = -1$. We need $\theta \le P_0(x)$, so $\theta \le \min_{\substack{x \in \{0, 1\}^n \\ n - \sum x_i \text{ is odd}}} P_0(x)$.

    Therefore, every possible joint PMF is of the form

    $$
    \Pr(X_1 = x_1, \dots, X_n = x_n) = \prod_{i=1}^n p_i^{x_i}(1-p_i)^{1-x_i} + \theta (-1)^{\sum_{i=1}^n (1-x_i)}
    $$

    and the admissible range of $\theta$ is

    $$
    \theta \in \left[- \min_{\sum x_i \equiv n \pmod 2} \prod_{i=1}^n p_i^{x_i}(1-p_i)^{1-x_i} , \min_{\sum x_i \not\equiv n \pmod 2} \prod_{i=1}^n p_i^{x_i}(1-p_i)^{1-x_i} \right]
    $$

    \item[(b)]
    
    $$
    \theta \in \left(-\frac{1}{16}, \frac{1}{16}\right)
    $$

    $$
    (-1)^{4 - \sum x_i} = (-1)^{4 - 2} = (-1)^2 = 1
    $$

    $$
    P(x) = \frac{1}{16} + \theta
    $$
    
    $$
    P(X_1 + X_2 + X_3 + X_4 = 2) = \sum_{\sum x_i = 2} P(x) = 6 \times \left(\frac{1}{16} + \theta\right) = \frac{3}{8} + 6\theta
    $$

    So $P(X_1 + X_2 + X_3 + X_4 = 2) \in \left(0, \frac{3}{4}\right)$.

    When $P(X_1 + X_2 + X_3 + X_4 = 2) = \frac{3}{8}$, we have $\theta = 0$. And we have $P(x) = P_0(x) + 0 = \prod_{i=1}^4 P(X_i = x_i)$, which is the condition for mutual independence. Therefore, the value $3/8$ does force mutual independence.
\end{enumerate}

\newpage
\begin{homeworkProblem}{20}
A population contains $m$ red balls and $m$ blue balls.
A uniformly chosen sample of size $n$ is drawn without replacement.
Let $X_{m,n}$ be the number of red balls, and let
$Y_n\sim\operatorname{Bin}(n,1/2)$.

\begin{enumerate}
\item[(a)] For fixed positive integer $n$ and fixed
$k\in\{0,\ldots,n\}$, give the PMF of $X_{m,n}$ and evaluate
\[
\lim_{m\to\infty}
m\left(\frac{\Pr(X_{m,n}=k)}{\Pr(Y_n=k)}-1\right).
\]
Do not merely quote convergence of the hypergeometric distribution.

\item[(b)] Set $n=2m-1$ and $A_m=\{m-1,m\}$. Evaluate
\[
\lim_{m\to\infty}
\left[\Pr(X_{m,2m-1}\in A_m)
-\Pr(Y_{2m-1}\in A_m)\right].
\]
Does a growing population with a fixed red proportion alone
justify the binomial approximation?
Prove the required limit without the central limit theorem
or Stirling's formula. 
\end{enumerate}
\end{homeworkProblem}

\solution

\begin{enumerate}
    \item[(a)]
    
    $X_{m,n} \sim \operatorname{HGeom}(2m, m, n)$

    $$
    \Pr(X_{m,n} = k) = \frac{\binom{m}{k} \binom{m}{n-k}}{\binom{2m}{n}}
    $$

    $$
    \Pr(Y_n = k) = \binom{n}{k} \left(\frac{1}{2}\right)^n
    $$

    $$
    \lim_{m \to \infty} m\left(\frac{\Pr(X_{m,n}=k)}{\Pr(Y_n=k)}-1\right)
    = \lim_{m \to \infty} m\left(\frac{\frac{\binom{m}{k} \binom{m}{n-k}}{\binom{2m}{n}}}{\binom{n}{k} \left(\frac{1}{2}\right)^n}-1\right)
    = \lim_{m \to \infty} m\left(\frac{\binom{m}{k} \binom{m}{n-k}}{\binom{2m}{n}\binom{n}{k}} 2^n - 1\right)
    $$

    $$
    \begin{aligned}
    \frac{\binom{m}{k} \binom{m}{n-k}}{\binom{2m}{n}\binom{n}{k}} 
    &= \frac{\frac{m!}{k!(m-k)!} \cdot \frac{m!}{(n-k)!(m-(n-k))!}}{\frac{(2m)!}{n!(2m-n)!} \cdot \frac{n!}{k!(n-k)!}} \\
    &= \frac{m! m! n! (2m-n)! k! (n-k)!}{k!(m-k)!(n-k)!(m-(n-k))!(2m)!n!} \\
    &= \frac{m! m! (2m-n)! }{(m-k)!(m-n+k)!(2m)!}
    \end{aligned}
    $$

    $$
    \frac{m!}{(m-k)!} = \prod_{i=0}^{k-1} (m-i)
    $$

    $$
    \frac{m!}{(m-n+k)!} = \prod_{j=0}^{n-k-1} (m-j)
    $$

    $$
    \frac{(2m-n)!}{(2m)!} = \frac{1}{\prod_{r=0}^{n-1} (2m-r)}
    $$

    $$
    \frac{\binom{m}{k} \binom{m}{n-k}}{\binom{2m}{n}\binom{n}{k}} 
    = \frac{\prod_{i=0}^{k-1} (m-i) \prod_{j=0}^{n-k-1} (m-j)}{\prod_{r=0}^{n-1} (2m-r)}
    $$

    $$
    \begin{aligned}
    \frac{\binom{m}{k} \binom{m}{n-k}}{\binom{2m}{n}\binom{n}{k}} 2^n
    &= \frac{\prod_{i=0}^{k-1} (m-i) \prod_{j=0}^{n-k-1} (m-j)}{\prod_{r=0}^{n-1} (m- \frac{r}{2})} \\
    &= \frac{\prod_{i=0}^{k-1} (1-\frac{i}{m}) \prod_{j=0}^{n-k-1} (1-\frac{j}{m})}{\prod_{r=0}^{n-1} (1- \frac{r}{2m})}
    \end{aligned}
    $$

    When $x \to 0$, we have $\prod (1 - a_i x) = 1 - x \sum a_i + O(x^2)$ and $\frac{1}{1 - cx} = 1 + cx + O(x^2)$.

    $$
    \prod_{i=0}^{k-1} \left(1 - \frac{i}{m}\right)
    = 1 - \frac{1}{m} \sum_{i=0}^{k-1} i + O\left(\frac{1}{m^2}\right)
    = 1 - \frac{k(k-1)}{2m} + O\left(\frac{1}{m^2}\right)
    $$

    $$
    \prod_{j=0}^{n-k-1} \left(1 - \frac{j}{m}\right)
    = 1 - \frac{1}{m} \sum_{j=0}^{n-k-1} j + O\left(\frac{1}{m^2}\right)
    = 1 - \frac{(n-k)(n-k-1)}{2m} + O\left(\frac{1}{m^2}\right)
    $$

    $$
    \prod_{r=0}^{n-1} \left(1 - \frac{r}{2m}\right)
    = 1 - \frac{1}{2m} \sum_{r=0}^{n-1} r + O\left(\frac{1}{m^2}\right)
    = 1 - \frac{n(n-1)}{4m} + O\left(\frac{1}{m^2}\right)
    $$

    $$
    \begin{aligned}
    \frac{\binom{m}{k} \binom{m}{n-k}}{\binom{2m}{n}\binom{n}{k}} 2^n
    &= \frac{\left(1 - \frac{k(k-1)}{2m} + O\left(\frac{1}{m^2}\right)\right)\left(1 - \frac{(n-k)(n-k-1)}{2m} + O\left(\frac{1}{m^2}\right)\right)}{1 - \frac{n(n-1)}{4m} + O\left(\frac{1}{m^2}\right)} \\
    &= \frac{1 - \frac{1}{m}\left[\frac{k(k-1)}{2} + \frac{(n-k)(n-k-1)}{2}\right] + O\left(\frac{1}{m^2}\right)}{1 - \frac{n(n-1)}{4m} + O\left(\frac{1}{m^2}\right)}
    \end{aligned}
    $$

    Use $\frac{1 - A/m}{1 - B/m} = 1 - \frac{A - B}{m} + O\left(\frac{1}{m^2}\right)$, we have

    $$
    \frac{\binom{m}{k} \binom{m}{n-k}}{\binom{2m}{n}\binom{n}{k}} 2^n
    = 1 - \frac{1}{m} \left( \frac{k(k-1) + (n-k)(n-k-1)}{2} - \frac{n(n-1)}{4} \right) + O\left(\frac{1}{m^2}\right)
    $$

    $$
    \frac{k(k-1) + (n-k)(n-k-1)}{2} - \frac{n(n-1)}{4}
    = \frac{n - (2k - n)^2}{4}
    $$

    $$
    \frac{\binom{m}{k} \binom{m}{n-k}}{\binom{2m}{n}\binom{n}{k}} 2^n - 1
    = \frac{n - (2k-n)^2}{4m} + O\left(\frac{1}{m^2}\right)
    $$

    $$
    \lim_{m \to \infty} m\left(\frac{\Pr(X_{m,n}=k)}{\Pr(Y_n=k)}-1\right)
    = \frac{n - (2k - n)^2}{4}
    $$

    \item[(b)]
    
    $$
    \Pr(X_{m,2m-1} \in A_m) = \Pr(X_{m,2m-1} = m-1) + \Pr(X_{m,2m-1} = m) = 1
    $$

    $$
    \Pr(Y_{2m-1} = m-1) = \binom{2m-1}{m-1} \left(\frac{1}{2}\right)^{2m-1}
    $$

    $$
    \Pr(Y_{2m-1} = m) = \binom{2m-1}{m} \left(\frac{1}{2}\right)^{2m-1} = \binom{2m-1}{m-1} \left(\frac{1}{2}\right)^{2m-1}
    $$

    $$
    \Pr(Y_{2m-1} \in A_m) = 2 \binom{2m-1}{m} \left(\frac{1}{2}\right)^{2m-1}
    = \binom{2m-1}{m} \left(\frac{1}{2}\right)^{2m-2}
    = \frac{\binom{2m}{m}}{4^m}
    $$

    Let $c_m = \frac{\binom{2m}{m}}{4^m} = \frac{(2m)!}{(m!)^2 4^m} = \prod_{k=1}^m \frac{2k-1}{2k}$.

    $$
    \left(\frac{2k-1}{2k}\right)^2 = \frac{2k-1}{2k} \cdot \frac{2k-1}{2k} < \frac{2k-1}{2k} \cdot \frac{2k}{2k+1} = \frac{2k-1}{2k+1}
    $$

    $$
    c_m^2 < \prod_{k=1}^m \frac{2k-1}{2k+1} = \frac{1}{3} \cdot \frac{3}{5} \cdots \frac{2m-1}{2m+1} = \frac{1}{2m+1}
    $$

    Hence $0 < c_m < \frac{1}{\sqrt{2m+1}}$.

    $$
    \lim_{m \to \infty} \Pr(Y_{2m-1} \in A_m) = \lim_{m \to \infty} c_m = 0
    $$

    Therefore, 

    $$
    \lim_{m \to \infty} \left[ \Pr(X_{m,2m-1} \in A_m) - \Pr(Y_{2m-1} \in A_m) \right] = 1 - 0 = 1
    $$
\end{enumerate}
\newpage
\begin{homeworkProblem}{20}
Let $X$ take values in $\{1,2,3,4\}$, with zero probabilities
allowed. Suppose
\[
\Pr(X\in\{1,2\})=\frac23,
\qquad
\Pr(X\in\{2,3\})=\frac12.
\]
Entropy is measured in bits, with $0\log_2 0=0$.

\begin{enumerate}
\item[(a)] Determine the maximum and minimum of $H(X)$ and
all PMFs attaining them. Prove that the extrema are global.

\item[(b)] If all four probabilities must be strictly positive,
determine the exact set of possible entropy values and state
which boundaries are attained. 
\end{enumerate}

\end{homeworkProblem}

\solution

\begin{enumerate}
    \item[(a)]
    
    Let $p_i = \Pr(X = i)$ for $i = 1, 2, 3, 4$, and $p_2 = c$.

    $$
    p_1 = \frac{2}{3} - c, \quad p_3 = \frac{1}{2} - c, \quad p_4 = 1 - p_1 - p_2 - p_3 = c - \frac{1}{6}
    $$

    $$
    c \in \left[\frac{1}{6}, \frac{1}{2}\right]
    $$

    $$
    \begin{aligned}
    H(c)
    &= -\sum_{i=1}^4 p_i \log_2 p_i \\
    &= -\left(\frac{2}{3} - c\right) \log_2 \left(\frac{2}{3} - c\right) - c \log_2 c - \left(\frac{1}{2} - c\right) \log_2 \left(\frac{1}{2} - c\right) - \left(c - \frac{1}{6}\right) \log_2 \left(c - \frac{1}{6}\right) \\
    &= -\frac{1}{\ln 2} \left[ \left(\frac{2}{3} - c\right) \ln \left(\frac{2}{3} - c\right) + c \ln c + \left(\frac{1}{2} - c\right) \ln \left(\frac{1}{2} - c\right) + \left(c - \frac{1}{6}\right) \ln \left(c - \frac{1}{6}\right) \right]
    \end{aligned}
    $$

    $$
    H'(c) = \frac{1}{\ln 2} \ln \left( \frac{\left(\frac{2}{3}-c\right)\left(\frac{1}{2}-c\right)}{c\left(c-\frac{1}{6}\right)} \right)
    $$

    $$
    \left(\frac{2}{3}-t\right)\left(\frac{1}{2}-t\right) = t\left(t-\frac{1}{6}\right)
    $$

    $$
    t = \frac{1}{3}
    $$

    $$
    H_{\max} = H\left(\frac{1}{3}\right) = \frac{1}{3} + \log_2 3, \quad \mathbf{p} = \left(\frac{1}{3}, \frac{1}{3}, \frac{1}{6}, \frac{1}{6}\right)
    $$
    
    $$
    H_{\min} = H\left(\frac{1}{6}\right) = H\left(\frac{1}{2}\right) = \frac{2}{3} + \frac{1}{2} \log_2 3
    ,\quad
    \mathbf{p} = \left(\frac{1}{2}, \frac{1}{6}, \frac{1}{3}, 0\right) \text{ or } \left(\frac{1}{6}, \frac{1}{2}, 0, \frac{1}{3}\right)
    $$

    \item[(b)]
    
    $$
    H \in \left(\frac{2}{3} + \frac{1}{2}\log_2 3, \; \frac{1}{3} + \log_2 3\right]
    $$
\end{enumerate}

\newpage
\begin{homeworkProblem}{20}
Let $N\ge2$ be an integer. A random variable $X$ takes values in
$\{0,1,\ldots,N\}$ and satisfies $\mathbb E[X]=\mu$, where
$0<\mu<N$. No other distributional assumptions are imposed.

\begin{enumerate}
\item[(a)] Determine the minimum and maximum possible variance
and characterize all equality cases. Cover both integer and
noninteger values of $\mu$. 
\item[(b)] Prove that every value between these extrema is
attainable, and provide an explicit PMF construction.

\end{enumerate}
\end{homeworkProblem}

\solution

\begin{enumerate}
    \item[(a)]
    
    $$
    \mathbb E[X] = \sum_{i=0}^N i p_i = \mu
    $$

    $$
    \mathbb E[X^2] = \sum_{i=0}^N i^2 p_i
    $$

    $$
    \Var[X] = \mathbb E[X^2] - (\mathbb E[X])^2 = \sum_{i=0}^N i^2 p_i - \mu^2
    $$

    When $\mu$ is an integer, we can construct $p_\mu = 1$, $p_i = 0$ for $i \ne \mu$, so $\Var[X] = 0$.

    When $\mu$ is not an integer, the minimum variance is $(\mu - \lfloor \mu \rfloor)(\lfloor \mu \rfloor + 1 - \mu)$ and the PMF is $p_{\lfloor \mu \rfloor} = \lfloor \mu \rfloor + 1 - \mu$, $p_{\lfloor \mu \rfloor + 1} = \mu - \lfloor \mu \rfloor$, $p_i = 0$ for $i \ne \lfloor \mu \rfloor, \lfloor \mu \rfloor + 1$.

    $$
    k(N - k) \ge 0 \implies k^2 \le N k.
    $$

    $$
    \mathbb{E}[X^2] \le N \mathbb{E}[X] = N \mu.
    $$

    $$
    \Var[X] = \mathbb{E}[X^2] - \mu^2 \le N \mu - \mu^2
    $$
    
    The equality holds if and only if $p_0 = 1 - \frac{\mu}{N}$ and $p_N = \frac{\mu}{N}$.

    \item[(b)]
    
    $$
    P_{\min}(k)
    =
    \begin{cases}
        1 - (\mu - \lfloor \mu \rfloor), & k = \lfloor \mu \rfloor \\
        \mu - \lfloor \mu \rfloor, & k = \lfloor \mu \rfloor + 1 \\
        0, & \text{otherwise}
    \end{cases}
    $$

    $$
    P_{\max}(k)
    =
    \begin{cases}
        1 - \frac{\mu}{N}, & k = 0 \\
        \frac{\mu}{N}, & k = N \\
        0, & \text{otherwise}
    \end{cases}
    $$
\end{enumerate}

\newpage
\begin{homeworkProblem}{20}
Consider independent Bernoulli trials with success probability
$p\in(0,1)$. Fix an integer $r\ge2$, and let $T_r$ be the
trial number of the $r$th success, including the successful trial.

Fix an integer $m\ge r-1$. An observer is told only that
$T_r>m$, not how many successes occurred during the first $m$
trials. On this event, let $R_m=T_r-m$.

\begin{enumerate}
\item[(a)] Find the complete conditional PMF of $R_m$ given
$T_r>m$. Express it also as a mixture of waiting-time
distributions, specifying the mixture weights.

\item[(b)] Can this conditional distribution be a first-success
distribution with any parameter? For each fixed integer $k\ge1$,
evaluate
\[
\lim_{m\to\infty}\Pr(R_m=k\mid T_r>m)
\]
and identify the limiting distribution.
\end{enumerate}
\end{homeworkProblem}

\solution

\begin{enumerate}
    \item[(a)]
    
    $T_r \sim \text{NBin}(r, p)$.

    $$
    P(T_r = k) = \binom{k-1}{r-1} p^r (1-p)^{k-r}, \quad k = r, r+1, \dots
    $$

    Define $S$ be the number of successes in the first $m$ trials. Then $S \sim \text{Bin}(m, p)$.

    $$
    P(S = k) = \binom{m}{k} p^k (1-p)^{m-k}, \quad k = 0, 1, \dots, m
    $$

    $T_r > m$, so $S < r$.

    $$
    \begin{aligned}
    P(S = k \mid T_r > m)
    &= P(S = k \mid S < r) \\
    &= \frac{P(S = k, S < r)}{P(S < r)} \\
    &=
    \begin{cases}
        \frac{P(S = k)}{P(S < r)}, & k = 0, 1, \dots, r-1 \\
        0, & k = r, r+1, \dots, m
    \end{cases}
    \\
    &=
    \begin{cases}
        \frac{\binom{m}{k} p^k (1-p)^{m-k}}{\sum_{i=0}^{r-1} \binom{m}{i} p^i (1-p)^{m-i}}, & k = 0, 1, \dots, r-1 \\
        0, & k = r, r+1, \dots, m
    \end{cases}
    \end{aligned}
    $$

    $$
    \begin{aligned}
    P(R_m = k \mid T_r > m)
    &= P(T_r - m = k \mid T_r > m) \\
    &= P(T_{r - S} = k \mid S < r) \\
    &= \sum_{s=0}^{r-1} P(T_{r - s} = k \mid S = s, S < r) P(S = s \mid S < r) \\
    &= \sum_{s=0}^{r-1} P(T_{r - s} = k) P(S = s \mid S < r) \\
    &=
    \sum_{s=0}^{r-1} \binom{k-1}{r-s-1} p^{r-s} (1-p)^{k-(r-s)} \frac{\binom{m}{s} p^s (1-p)^{m-s}}{\sum_{i=0}^{r-1} \binom{m}{i} p^i (1-p)^{m-i}}, & k = 1, 2, \dots \\
    &=
    \frac{p^r (1-p)^{k+m-r}}{\sum_{i=0}^{r-1} \binom{m}{i} p^i (1-p)^{m-i}} \sum_{s=0}^{r-1} \binom{k-1}{r-s-1} \binom{m}{s} , & k = 1, 2, \dots \\
    &=
    \frac{\binom{k+m-1}{r-1} p^r (1-p)^{k+m-r}}{\sum_{i=0}^{r-1} \binom{m}{i} p^i (1-p)^{m-i}}, & k = 1, 2, \dots 
    \end{aligned}
    $$

    $$
    P(R_m = k \mid T_r > m) = \frac{P(T_r = m+k)}{P(T_r > m)} = \frac{\binom{m+k-1}{r-1} p^r (1-p)^{m+k-r}}{\sum_{i=0}^{r-1} \binom{m}{i} p^i (1-p)^{m-i}}
    $$

    $$
    P(R_m = k \mid T_r > m) = \sum_{s=0}^{r-1} w_s \cdot P(W_{r-s} = k)
    $$

    and $w_s = P(S = s \mid S < r) = \frac{\binom{m}{s} p^s (1-p)^{m-s}}{\sum_{i=0}^{r-1} \binom{m}{i} p^i (1-p)^{m-i}}, \quad s \in \{0, 1, \dots, r-1\}$.
    
    \item[(b)]
    
    $$
    \frac{P(R_m = k+1 \mid T_r > m)}{P(R_m = k \mid T_r > m)} = \frac{\binom{m+k}{r-1} (1-p)}{\binom{m+k-1}{r-1}} = (1-p) \frac{m+k}{m+k-r+1} = (1-p) \left( 1 + \frac{r-1}{m+k-r+1} \right)
    $$

    It cannot be a first-success distribution.

    $$
    \frac{\binom{m+k-1}{r-1}}{\binom{m}{r-1}} p^r (1-p)^{k-r} \xrightarrow{m \to \infty} 1 \cdot p^r (1-p)^{k-r} = p^r (1-p)^{k-r}
    $$

    $$
    \sum_{i=0}^{r-1} \frac{\binom{m}{i}}{\binom{m}{r-1}} p^i (1-p)^{-i} \xrightarrow{m \to \infty} p^{r-1} (1-p)^{-(r-1)}
    $$

    $$
    \lim_{m \to \infty} P(R_m = k \mid T_r > m) = \frac{p^r (1-p)^{k-r}}{p^{r-1} (1-p)^{-(r-1)}} = p(1-p)^{k-1}, \quad k = 1, 2, \dots
    $$
\end{enumerate}
\end{document}